\documentclass[sigconf]{acmart}
\settopmatter{printacmref=false}
\setcopyright{none}
\title{FounderBench Supplementary Appendix}
\author{Yufeng Wang}
\affiliation{%
  \institution{Independent Researcher}
  \country{United States}}
\email{louiswang524@gmail.com}
\renewcommand{\shortauthors}{Yufeng Wang}

\begin{document}
\maketitle

\section{Task Families}
FounderBench v0.3.0 contains five fixed tasks in each family.
Market selection rewards sufficient research, viable-market choice, and preserved cash.
First revenue rewards customers, recurring revenue, runway, and reputation.
Retention improvement rewards quality, support, customer preservation, and focus.
Churn recovery rewards stabilized customers, reduced churn, restored reputation, and adequate capacity.
Demo Day traction rewards growth, recurring revenue, cash, and credibility.
Pricing rewards movement toward sustainable market-relative prices.
Runway preservation rewards survival and disciplined cost reduction.
Pivot decision rewards diagnosis and timely market change.
Fundraising rewards credible traction before simulated capital raising.
Channel expansion rewards scaling working offers through campaigns and partnerships.

Complete human-readable cards and executable setup and scoring functions are included in the artifact.

\section{Prompt and Response Contract}
Each provider receives the task description, current structured observation, allowed action schemas, and recent memory.
The expected response root is:
\begin{verbatim}
{
  "rationale": "brief explanation",
  "actions": [
    {"type": "research_market", ...}
  ]
}
\end{verbatim}
Only entries in \texttt{actions} are executable.
The parser extracts a JSON object, validates the root and action list, checks required fields and numeric types, and records a diagnostic category when parsing fails.
The exact generated protocol and prompt hashes are released in
\texttt{outputs/founderbench-prompt-protocol.md}.

\section{Core Transition Sketch}
After actions execute, each offer settles weekly customers with seeded Poisson draws.
With awareness $a$, quality $q$, price $p$, reputation $r$, market demand $d$, competition $c$, willingness to pay $w$, and support load $\ell$:
\begin{align*}
\mathrm{price\_fit}&=\max(0.08,1.25-p/\max(1,w)),\\
\lambda_{\mathrm{new}}&=\max\bigl(0,(d\,a\,q\,\mathrm{price\_fit}\,r-0.18c)\,5.5\bigr),\\
\rho&=\max\bigl(0.01,0.08\ell+(0.65-q)\,0.05-0.025r\bigr).
\end{align*}
Revenue is $p\cdot\mathrm{customers}$; support cost scales with $\ell\cdot\mathrm{customers}$.
Support overload reduces reputation when required support exceeds capacity.
Research replaces noisy signals with exact market values; campaigns raise awareness; improve\_offer raises quality; hire\_agent raises capacity; cut\_cost reduces burn; raise\_funding injects cash under traction/reputation gates.
Market momentum evolves with a seeded volatility shock.
All constants and action branches are in the released \texttt{env.py}.

\section{Score Construction}
Every task score is a sum of documented positive components and penalties, clamped to $[0,100]$.
Market selection:
\begin{multline*}
f=25\min(1,R/3)+55\cdot\mathbf{1}[\mathrm{good\_market}]\\
\quad +20\min(1,\mathrm{cash}/8500).
\end{multline*}
First revenue (targets $C^\star,v^\star,K^\star$ vary by variant):
\begin{multline*}
f=45\min(1,C/C^\star)+25\min(1,v/v^\star)\\
\quad +20\min(1,\mathrm{cash}/K^\star)+10\min(1,r/0.65).
\end{multline*}
Other families follow the same weighted-component pattern (quality/reputation/churn for retention; capacity for churn shock; growth for Demo Day; price-band fit for pricing; survival for runway; pivot correctness for pivot; funding gates for fundraising; scale for channel).
The complete per-family rubric is in \texttt{outputs/founderbench-score-rubric.md}.

The official leaderboard metric is the mean task score; we also report median and IQR across tasks.
A task is solved at score $\geq70$, chosen as a fixed ``clearly successful'' operating cut between the generic-heuristic and task-aware means rather than a tuned optimum; we do not sweep alternative thresholds here.
``Pub.\ test'' summarizes FND-031--050; both splits are visible.
Task bootstrap intervals resample the 50 score observations and estimate sensitivity to task composition, not provider sampling.

\section{Task-Aware Heuristic Summary}
\texttt{TaskHeuristicPolicy} first maps task ID to family index $({\mathrm{id}}-1)//5$, then applies a short family branch; otherwise it falls back to the generic heuristic.
Representative rules:
\begin{itemize}
  \item Market selection / first revenue: research top unresearched markets early, then \texttt{build\_offer} on the best demand--price--competition score.
  \item Retention: improve quality and support when below thresholds.
  \item Churn shock: hire if capacity is low, then support/improve.
  \item Demo Day: improve, campaign, hire under load, support.
  \item Pricing: interview then set price near $0.78w$, optionally campaign.
  \item Runway: \texttt{cut\_cost}, light support, selective improve.
  \item Pivot: research target market; pivot when stalled; reprice/improve/campaign.
  \item Fundraising: support/improve then \texttt{raise\_funding}.
  \item Channel: \texttt{partner\_channel}, improve, campaign, hire under load.
\end{itemize}
This policy uses family identity unavailable to ordinary hosted prompts unless leaked by the task text; treat it as a calibration ceiling, not a human expert.
On the frozen private remix, deterministic calibration alone already reverses the public ladder (task-aware below generic heuristic), which is a warning against reading the public 80.90 gap as stable strategic headroom.

\section{Diagnostics}
Outputs include invalid actions, over-budget decisions, provider-error categories, action counts, simulated API cost, latency, and token counts when available.
Audit mode additionally records model identifier, task and week, prompt hash, redacted response, usage, estimated cost, and call latency.
Automated redaction is a safety layer and traces must be inspected before public release.

Provider/parser failures use documented fallback behavior and remain in official scores.
The paper's sensitivity artifact reports affected-task and unaffected-task means without treating them as a corrected counterfactual.
\begin{table*}[t]
\caption{Provider-error sensitivity with normalized rates. Err/task is error decisions divided by 50 tasks; Aff.\ tasks is the count (and fraction) of tasks with $\geq$1 error decision. Means stratify observed tasks and are not counterfactual corrected scores.}
\label{tab:errors}
\centering\scriptsize
\begin{tabular}{lrrrrrl}
\toprule
Model & Errs & Err/task & Aff.\ tasks & Aff.\ mean & Clean mean & Dominant category \\
\midrule
Gemini 3.5 Flash & 59 & 1.18 & 25 (50\%) & 73.57 & 61.80 & provider\_rate\_limit (30) \\
Grok 4.5 & 0 & 0.00 & 0 (0\%) & -- & 66.53 & -- \\
GPT-5.6 Sol & 0 & 0.00 & 0 (0\%) & -- & 66.39 & -- \\
Kimi K3 & 70 & 1.40 & 11 (22\%) & 46.47 & 71.04 & provider\_rate\_limit (64) \\
Claude Sonnet 5 & 0 & 0.00 & 0 (0\%) & -- & 63.90 & -- \\
DeepSeek V4 Reasoner & 3 & 0.06 & 3 (6\%) & 60.66 & 62.55 & invalid\_json (2) \\
Claude Sonnet 4.5 & 0 & 0.00 & 0 (0\%) & -- & 61.09 & -- \\
DeepSeek Chat & 0 & 0.00 & 0 (0\%) & -- & 56.59 & -- \\
GLM 4.5 Air & 0 & 0.00 & 0 (0\%) & -- & 54.78 & -- \\
Gemini 2.5 Flash & 340 & 6.80 & 44 (88\%) & 50.33 & 70.03 & provider\_io\_error (332) \\
Grok 4.3 & 3 & 0.06 & 3 (6\%) & 44.23 & 53.12 & no\_json\_object (2) \\
\bottomrule
\end{tabular}
\end{table*}


\section{Submission Validation}
The validator requires:
\begin{itemize}
  \item exactly 50 unique official task identifiers;
  \item score bounds and pass flags consistent with the threshold;
  \item aggregate solved, solve-rate, and average-score consistency;
  \item public development/test summaries;
  \item execution diagnostics and declared policy/model metadata; and
  \item complete result objects for every task.
\end{itemize}
The frozen paper registry records each accepted evidence path, submission report, benchmark version, model label, run seed, headline outcomes, and SHA-256 digest.
For newer audited runs it also records call-level model identifiers.
Some earlier validated rows predate that audit field; their model labels are frozen from the variant-specific execution configuration and evidence path.
Future submission schemas should require the exact API model identifier at the run root.

\section{Artifact Map}
The paper package contains:
\begin{itemize}
  \item the main manuscript, generated tables, and publication figures;
  \item table, figure, and result-analysis generation scripts;
  \item \texttt{outputs/founderbench-paper-model-registry.json};
  \item \texttt{outputs/founderbench-paper-analysis.json};
  \item complete validated task-level provider JSON and submission reports; and
  \item benchmark card, datasheet, responsible-use report, reproduction guide,
        environment report, and integrity manifests.
\end{itemize}

\section{Additional Limitations}
The environment excludes many real organizational concerns, including law, accounting, hiring, workplace safety, data governance, and interpersonal dynamics.
Actions are coarse abstractions.
The synthetic markets encode designer assumptions and may admit simulator-specific strategies.
The five tasks per family provide diagnostic breadth but limited statistical resolution.
Provider model aliases can change after collection, and the single-run study does not quantify inference variance.
No human-subject evaluation, founder comparison, official private leaderboard, or real-world outcome validation is claimed.
Recommended upgrades: repeated hosted runs, standardized retry budgets, stochastic task variants, expert calibration, and a versioned private split.

\end{document}
